\section{All-time upper covers and theorem assembly}
\label{hd:cover-section}

We now construct finite covers of the whole trace in each bounded box.
The two charges are different: disjoint successful balls are paid for
by the same actual terminal mass, while missed cells are paid for by
the cleanup entry. The finite changes of law estimate adapted test
failures; the cover itself is selected afterward.

The use of regions with large trace mass together with a finer residual
cover already appears in Rezaei
\cite[Section~3, equations~(9)--(11)]{Rezaei2018}; see also the
conditional covering sketch in
\cite[Remark~4.5]{Zhan2019Optimal}.
Here the all-moment batches choose the thresholds and scales, and finite
adapted tests give the conditional failure bound. We then account for
every missed cell, including cells near the real boundary.

We include cells meeting the real axis. Their centers will always
lie strictly in \(\HH\), even though the cells themselves may
contain boundary points.

\begin{lemma}[A cell visit forces a conformal-radius crossing]
\label{hd:visit}
Let \(Q\) be a closed square of side \(r\), centered at
\(z=x+iy\), where \(y>0\).
If \(y>4r\) and \(Q\cap\Gamma\ne\varnothing\), then \(z\)
reaches conformal radius \(8r\) while alive.
For every \(r>0\) and \(y>0\),
\begin{equation}
 \PP(Q\cap\Gamma\ne\varnothing)
 \le C_a r^\alpha y^{-\alpha}.
 \label{hd:one-cell}
\end{equation}
\end{lemma}

\begin{proof}
Suppose the radius never reaches \(8r\) before the lifetime
\(T_z\) of \(z\). Since \(R_0(z)=2y>8r\), continuity and Koebe's
Quarter Theorem give
\[
 B(z,2r)\subset H_t\qquad(t<T_z).
\]
The past trace avoids this open ball. If \(T_z<\infty\),
continuity prevents the trace from entering its interior at
time \(T_z\) as well.
The ball lies in \(\HH\), is connected, and avoids
\(\gamma([0,T_z])\). It therefore lies in the same complementary
component as \(z\). This component is bounded, since \(z\) has
been swallowed. Each point of the ball has an open neighborhood
in that component, disjoint from \(H_{T_z}\); hence
\[
 B(z,2r)\cap\overline{H_{T_z}}=\varnothing.
\]
The future trace lies in \(\overline{H_{T_z}}\), so it cannot
enter the ball. If \(T_z=\infty\), the ball is avoided at every
finite time directly. In both cases \(Q\), which lies strictly
inside this ball, is unvisited.
This proves the crossing assertion, including when
\(4<\kappa<8\).

If \(y/r\) exceeds a sufficiently large constant depending on
\(a\), the inner time to radius \(8r\) is
\(\sigma=(2a)^{-1}\log(y/(4r))\), long enough for
Lemma~\ref{op:mixing}. At its hitting time \(\tau\),
Lemma~\ref{op:finite-tilt} gives
\[
 \PP(\tau<\infty)
 =(8r)^\alpha G_0(z)\,
     \EE_z^*[\sin^{-p}\theta_\tau(z)]
 \le C_a r^\alpha y^{-\alpha}.
\]
Here \(G_0(z)\le(2y)^{-\alpha}\), and
\[
 \int_0^\pi\sin^{-p}\theta\,\pi_a(\dd\theta)
 =c_a\int_0^\pi\sin\theta\,\dd\theta<\infty.
\]
For the remaining values of \(y/r\), a larger \(C_a\) makes the
right side of Equation~\eqref{hd:one-cell} at least one, so the
trivial probability bound completes the proof.
\end{proof}

\begin{proposition}[The upper Hausdorff bound]
\label{hd:upper}
For every integer $m\ge1$, the random variable
$\Haus^h(\Gamma\cap B_m)$, where $B_m=[-m,m]+i[0,m]$, is measurable
and satisfies
\begin{equation}
 \EE\,\Haus^h(\Gamma\cap B_m)<\infty.
 \label{hd:expected-box}
\end{equation}
In particular, these whole-trace bounded-box measures are almost
surely finite simultaneously for all $m$.
\end{proposition}

\begin{proof}
Fix $m$. We first estimate finite adapted test failures at the final
mesh of a batch. Then we choose a finite terminal-mass cover and bound
its expected cost. The last step establishes measurability and passes
to arbitrarily small diameters.

\paragraph{Finite tests at the final mesh.}
Take a sufficiently large batch $n$ and write
\[
 J=I_n-1,\qquad r_{\mathrm f}=r_{n,I_n},\qquad
 r_i=r_{ni},\quad b_i=b_{ni}\quad(1\le i\le J).
\]
There is at least one test for all sufficiently large \(n\), by
Lemma~\ref{gauge:multiplicity}. Cover the box
\(B_m=[-m,m]+i[0,m]\) by finitely many closed squares of side
\(r_{\mathrm f}\), beginning at \(x=-m\) and at height zero,
and retain the whole squares in the outermost rows and columns.
Their centers have heights
\[
 y_j=(j+\tfrac12)r_{\mathrm f},\qquad j\ge0.
\]
At every grid center \(z\), the closed ball
\(\overline B(z,C_*r_i)\) is a \emph{candidate} for entry \(i\)
if
\begin{equation}
 \mu_\infty\bigl(\overline B(z,C_*r_i)\cap\HH\bigr)
       \ge c_*^d b_i r_i^d.
 \label{hd:candidate}
\end{equation}
Every such ball contains the final cell centered at \(z\).
There are only finitely many candidate checks in a batch.
They involve terminal mass; candidate selection itself is made
after the sample has been realized.

We estimate the probability that a cell is visited and has no
candidate through finite stopped tests. Consider a center \(z\) with \(y\ge e^{-n}\).
For large \(n\), Lemma~\ref{hd:visit} shows that a visit requires
the finite hitting time \(\tau_{\mathrm f}\) of radius
\(8r_{\mathrm f}\). For each \(i\), let \(S_i\) and \(T_i\) be
the first times, while \(z\) remains alive, that its conformal radius
reaches \(100r_i\) and \(100r_i e^{-2a\ell_n}\), respectively;
set either time to infinity on failure. On
\(\{\tau_{\mathrm f}<\infty\}\), both times are finite and the
inner-clock interval from \(S_i\) to \(T_i\) has length \(\ell_n\).
On \(\{T_i<\infty\}\), let \(F_i:\DD\to H_{S_i}\) be the disk map used in
Equation~\eqref{hd:spatial-comparison}, and define the stopped test mass
\begin{equation}
 X_i=\int_{F_i(K)}
       |F_i'(F_i^{-1}(x))|^{-d}
       (\mu_{T_i}-\mu_{S_i})(\dd x).
 \label{hd:stopped-test}
\end{equation}
Set \(X_i=0\) on \(\{T_i=\infty\}\).
The variable is measurable at \(T_i\) and, on a reached endpoint,
is the canonical disk mass \(L_{\ell_n}\) for the restart angle.
Only paths reaching
\(\tau_{\mathrm f}\) will enter the stopped event below.
Success, \(X_i\ge b_i\), implies
\[
 (\mu_{T_i}-\mu_{S_i})(F_i(K))\ge c_*^d b_i r_i^d,
\]
so Equation~\eqref{hd:koebe-sets} and monotonicity of mass force
the ball \(\overline B(z,C_*r_i)\) to qualify in
Equation~\eqref{hd:candidate}.

The gaps in the inner clock are deterministic;
Figure~\ref{fig:batch-timeline} shows their order.
Writing \(\sigma_i=s_z(S_i)\) and
\(\sigma_{\mathrm f}=s_z(\tau_{\mathrm f})\), the radius formula
gives exactly
\begin{align}
 \sigma_1
   &=\frac{1}{2a}\log\frac{2y}{100e^{-n^2}}
     \ \ge\ \frac{n^2-n-\log50}{2a}, \notag\\
 \sigma_{i+1}-(\sigma_i+\ell_n)&=n
       \qquad(1\le i<J), \label{hd:test-gaps}\\
 \sigma_{\mathrm f}-(\sigma_J+\ell_n)
   &=n+\frac{1}{2a}\log\frac{100}{8}. \notag
\end{align}
\begin{figure}[t]
\centering
\begin{tikzpicture}[x=1cm,y=1cm,>=Latex,font=\small]
  \draw[->,thick] (0,0)--(13,0) node[right] {\(s\)};
  \foreach \x/\y in {0.8/2.6,3.8/5.6,7.3/9.1} {
    \fill[blue!9] (\x,0.15) rectangle (\y,0.82);
    \draw[blue!60!black] (\x,0.15) rectangle (\y,0.82);
    \node at ({(\x+\y)/2},0.49) {test: \(\ell_n\)};
    \draw (\x,0.07)--(\x,-0.07);
    \draw (\y,0.07)--(\y,-0.07);
  }
  \node[below] at (0.8,-0.1) {\(\sigma_1\)};
  \node[below] at (2.6,-0.1) {\(\sigma_1+\ell_n\)};
  \node[below] at (3.8,-0.1) {\(\sigma_2\)};
  \node[below] at (5.6,-0.1) {\(\sigma_2+\ell_n\)};
  \node[below] at (7.3,-0.1) {\(\sigma_J\)};
  \node[below] at (9.1,-0.1) {\(\sigma_J+\ell_n\)};
  \node at (6.45,0.48) {\(\cdots\)};
  \draw[<->] (2.6,1.1)--(3.8,1.1)
    node[midway,above] {\(n\)};
  \draw[<->] (9.1,1.1)--(12.2,1.1)
    node[midway,above] {\(n+\frac{1}{2a}\log\frac{100}{8}\)};
  \draw[dashed] (12.2,-0.07)--(12.2,0.85);
  \node[below] at (12.2,-0.1) {\(\sigma_{\mathrm f}\)};
  \node[align=center,font=\footnotesize] at (6.5,-1.15)
    {All failures are known at \(T_J\). The final gap controls\\
     the inverse-angle factor in the change back to the original law.};
\end{tikzpicture}
\caption{One batch at a fixed grid center, in inner time (schematic).
There are \(J=I_n-1\) mass tests, each of length \(\ell_n\),
with gaps \(n\) between consecutive tests. The coordinates are
\(s_z(S_i)=\sigma_i\), \(s_z(T_i)=\sigma_i+\ell_n\), and
\(s_z(\tau_{\mathrm f})=\sigma_{\mathrm f}\); the final coordinate
corresponds to conformal radius \(8r_{n,I_n}\).
The first restart is preceded by a mixing gap as well.
The tests are conditioned on their actual past; independence is unnecessary.}
\label{fig:batch-timeline}
\end{figure}

For the failure estimate, we now work under the finite centered law
stopped at \(\tau_{\mathrm f}\), under which all these levels are reached.
For all sufficiently large \(n\), each of these gaps permits
the uniform mixing estimate. Conditional on the entire stopped
history through \(T_{i-1}\), the angle at \(S_i\) has density at
least \(1/2\) relative to \(\pi_a\); the first test has the same
bound from its initial gap.
Conditional on the history through \(S_i\), the disk law of
\(X_i\) depends only on that angle, by
Proposition~\ref{mass:restart} and the finite conditional tilt.
Thus, with
\[
 q_i=\PP^*_{\mathrm{mix}}(L_{\ell_n}\ge b_i),
\]
the conditional probability of success is at least \(q_i/2\).
Iterated conditional expectations, followed by
Lemma~\ref{gauge:multiplicity}, give
\begin{equation}
 \PP_z^*(X_i<b_i\text{ for every }i\le J)
 \le\prod_{i=1}^J(1-q_i/2)
 \le \exp\left(-\tfrac12\sum_{i=1}^Jq_i\right)
 \le \exp(-2^n/16).
 \label{hd:test-failure}
\end{equation}
This argument conditions on the actual past at each test.
It requires no independence between tests or between centers.

\paragraph{Changing law for the failure event.}
The event that a cell is visited and has no candidate is
contained in the larger event
\begin{equation}
 A=\{\tau_{\mathrm f}<\infty\}
       \cap\bigcap_{i=1}^J\{X_i<b_i\}.
 \label{hd:adapted-enlargement}
\end{equation}
All test endpoints precede \(\tau_{\mathrm f}\), so \(A\) is
measurable at that stopping time. Only this enlarged event is
used in changing law. In particular, no terminal candidate
information is put into a finite Radon--Nikodym formula.
Let \(\widehat\theta=\theta_{\tau_{\mathrm f}}(z)\) and
\(F_{\mathrm{fail}}=\bigcap_i\{X_i<b_i\}\).
The final gap in Equation~\eqref{hd:test-gaps}, conditional on
the history through \(T_J\), gives
\[
 \EE_z^*[\sin^{-p}\widehat\theta\mid\mathcal F_{T_J}]
 \le 2\int\sin^{-p}\theta\,\pi_a(\dd\theta).
\]
Since \(F_{\mathrm{fail}}\) is known at \(T_J\), the density at
\(\tau_{\mathrm f}\) and Equation~\eqref{hd:test-failure} yield
\begin{align}
 &\PP(\text{the cell at }z\text{ is visited and has no candidate})
 \notag\\
 &\quad\le\PP(A)
  =(8r_{\mathrm f})^\alpha G_0(z)\,
       \EE_z^*[\sin^{-p}\widehat\theta;\ F_{\mathrm{fail}}]
 \notag\\
 &\quad\le C_a r_{\mathrm f}^{\alpha}y^{-\alpha}
                       \exp(-2^n/16)
       \qquad(y\ge e^{-n}).
 \label{hd:upper-cell}
\end{align}
For centers below \(e^{-n}\), use the unconditional estimate
\begin{equation}
 \PP(\text{the cell at }z\text{ is visited})
 \le C_a r_{\mathrm f}^{\alpha}y^{-\alpha}
 \label{hd:lower-cell}
\end{equation}
from Lemma~\ref{hd:visit}. This includes the bottom row:
its center has height \(r_{\mathrm f}/2\), and the trivial
probability bound is absorbed by the displayed estimate.

\paragraph{Selecting and charging the cover.}
Order all candidate balls by
decreasing radius, using a fixed deterministic order for ties.
Choose the first ball, delete every ball meeting it, and
continue until no candidates remain. The selected closed balls
are pairwise disjoint. A deleted ball has radius no greater
than that of the ball that deletes it, so it is contained in
the concentric triple of that selected ball.
The triples therefore cover every final cell having a candidate.
For a chosen ball associated with \((b_i,r_i)\), its triple
has diameter \(6C_*r_i\); Equation~\eqref{gauge:definition}
and its qualification give
\begin{align}
 h(6C_*r_i)
 &\le (6C_*)^2 b_i r_i^d \notag\\
 &\le \frac{(6C_*)^2}{c_*^d}\,
       \mu_\infty\bigl(\overline B(z,C_*r_i)\cap\HH\bigr).
 \label{hd:candidate-cost}
\end{align}
For large \(n\) all these balls lie, after intersection with
\(\HH\), in the fixed bounded set
\[
 \widetilde B_m=[-m-2,m+2]+i(0,m+2].
\]
Disjointness thus bounds their total gauge cost by
\[
 \frac{(6C_*)^2}{c_*^d}\,\mu_\infty(\widetilde B_m).
\]
This is the actual mass of a fixed region on the same sample.
It has finite expectation, even at the real boundary:
\begin{equation}
 \EE\mu_\infty(\widetilde B_m)
 =\int_{\widetilde B_m}G_0(z)\,\dd A(z)
 \le \frac{2^{-\alpha}(2m+4)(m+2)^{1-\alpha}}{1-\alpha}
 <\infty.
 \label{hd:boundary-integrability}
\end{equation}
The mean identity on this set follows by an interior compact
exhaustion and monotone convergence.

\paragraph{Paying for missed cells and the boundary.}
Add every final cell that meets $\Gamma$ and has no candidate.
The cleanup entry gives the individual cost
\[
 h(\sqrt2\,r_{\mathrm f})\le2n r_{\mathrm f}^d.
\]
There are \(O_m(r_{\mathrm f}^{-1})\) columns. Since
\(0<\alpha<1\), summing the half-integer row heights gives
\begin{align*}
 \sum_{\text{all cells}}y^{-\alpha}
     &\le C_{a,m}r_{\mathrm f}^{-2},\\
 \sum_{y<e^{-n}}y^{-\alpha}
     &\le C_{a,m}r_{\mathrm f}^{-2}
                        (e^{-n}+r_{\mathrm f})^{1-\alpha}.
\end{align*}
For example,
\(\sum_{j=0}^N(j+\tfrac12)^{-\alpha}
\le C_\alpha(N+1)^{1-\alpha}\) proves both bounds.
Multiplying Equations~\eqref{hd:upper-cell} and
\eqref{hd:lower-cell} by the cleanup cost, and using
\(d+\alpha=2\), bounds the expected added cost by
\begin{equation}
 C_{a,m}\left[
       n e^{-2^n/16}
       +n(e^{-n}+r_{\mathrm f})^{1-\alpha}
       \right].
 \label{hd:cleanup-cost}
\end{equation}
It is bounded uniformly in \(n\) and tends to zero.
In particular, the bottom row contributes at most
\(C_m n r_{\mathrm f}^{d-1}\). Thus points of the trace on
the real axis are included in the cover, although the auxiliary
mass itself is defined on \(\HH\).

\paragraph{Measurability and the limiting cover.}
These constructions are measurable. The candidate family is
finite with measurable terminal-mass indicators, and the
greedy order is deterministic. For a fixed closed cell, the
event of meeting the trace is measurable on every compact time
interval, by continuity; a countable union over integer time
horizons gives the full trace event.
All upper-bound centers, entries, and boxes constitute a
countable family, so the stopped measure identities used in
constructing their tests can be imposed on one event.

Let \(C_n\) be the total gauge cost of the resulting cover of
\(\Gamma\cap B_m\). Equations~\eqref{hd:candidate-cost}--%
\eqref{hd:cleanup-cost} show that
\(\sup_n\EE C_n<\infty\), after omitting finitely many initial
batches. Every cover diameter is at most
\[
 \delta_n=6C_*e^{-n^2},
\]
which tends deterministically to zero.
Fatou's Lemma gives
\[
 \EE[\liminf_{n\to\infty}C_n]
 \le\liminf_{n\to\infty}\EE C_n<\infty,
\]
so the liminf cost is finite almost surely.
For every fixed \(\delta>0\), all sufficiently late covers are
admissible for \(\Haus^h_\delta(\Gamma\cap B_m)\), and hence
\[
 \Haus^h_\delta(\Gamma\cap B_m)
 \le\liminf_{n\to\infty}C_n.
\]
Letting $\delta\downarrow0$ gives the pathwise domination
\begin{equation}
 \Haus^h(\Gamma\cap B_m)\le\liminf_{n\to\infty}C_n.
 \label{hd:liminf-domination}
\end{equation}

We justify ordinary expectation of the left side. For an integer $T$,
$K_T=\gamma([0,T])\cap B_m$ is a measurable compact random set.
Indeed, for a fixed closed $F$, the event $K_T\cap F\ne\varnothing$
is the measurable event
$\min_{0\le t\le T}\dist(\gamma(t),B_m\cap F)=0$; the minimum can
be computed on rational times. Empty target sets give the empty event.
Hit events for open sets follow by an increasing union of closed
subsets, so these tests give the required compact-set measurability.
Intersection with $B_m$ need not depend continuously on the path.

For completeness, $K\mapsto\Haus^h(K)$ is measurable on compact sets.
Use the countable family of finite unions of open balls with rational
centers and radii. The limit of the corresponding countable infima
over finite covers, with their diameter cutoff tending to zero, equals
$\Haus^h(K)$. To see this, start with any countable cover of the compact
$K$ at a smaller cutoff. Enlarge its sets to open neighborhoods with
arbitrarily small diameter slack and summable cost slack, using
continuity of $h$ (also at zero). Compactness gives a finite subcover.
Shrink that finite cover on $K$ to compact pieces, and cover each piece
by finitely many rational balls inside its assigned neighborhood.
The grouped balls remain inside that neighborhood, preserving its
diameter bound and cost. This proves equality after letting the cutoff
and slack tend to zero; equality at an unchanged fixed cutoff is not
needed. Each admissibility event $K\subset O$ for a fixed open $O$
is measurable, so the countable infima and their limit are measurable.

Finally $\Gamma\cap B_m=\bigcup_{T\in\mathbb N}K_T$ is an increasing
union of compact sets. Hausdorff outer measure is a measure on Borel
sets, so continuity from below gives
$\Haus^h(\Gamma\cap B_m)=\lim_T\Haus^h(K_T)$.
This is measurable. Taking expectations in
\eqref{hd:liminf-domination} and applying the preceding Fatou bound
proves \eqref{hd:expected-box}. A countable intersection now gives
almost-sure finiteness in every integer box.
\end{proof}

\begin{proof}[Proof of Theorem~\ref{thm:main}]
Equation~\eqref{coef:formula} specifies the coefficients
deterministically from \(\kappa\), and
Proposition~\ref{gauge:admissible} shows that
Equation~\eqref{gauge:definition} defines an admissible gauge.
Intersect the probability-one events of
Propositions~\ref{hd:lower} and~\ref{hd:upper}.
On this one event, the lower bound holds for every real
\(0<s<t<\infty\).
The image of each finite time interval is bounded by continuity
of \(\gamma\), and therefore lies in one of the integer boxes
of Proposition~\ref{hd:upper}; its Hausdorff measure is finite.
Thus the required positive finite measure holds simultaneously
for every such interval. The constant change of capacity time
used at the start preserves this all-interval assertion.
\end{proof}
